\section{AMC 8 2011}
        \begin{enumerate}
        	\item (\textit{AMC 8 2011}) Margie bought $3$ apples at a cost of $50$ cents per apple. She paid with a 5-dollar bill. How much change did Margie receive?


 $\textbf{(A) }\ \textdollar 1.50 \qquad \textbf{(B) }\ \textdollar 2.00 \qquad \textbf{(C) }\ \textdollar 2.50 \qquad \textbf{(D) }\ \textdollar 3.00 \qquad \textbf{(E) }\ \textdollar 3.50$

 \textbf{Answer: }E



\textbf{Solution}

$50$ cents is equivalent to $\textdollar 0.50.$ Then the three apples cost $3 \times \textdollar 0.50 = \textdollar 1.50.$ The change Margie receives is $\textdollar 5.00 - \textdollar 1.50 = \boxed{\textbf{(E)}\ \textdollar 3.50}$

 -Mathloveraldoxxx5223

 


	\item (\textit{AMC 8 2011}) Karl's rectangular vegetable garden is $20$ feet by $45$ feet, and Makenna's is $25$ feet by $40$ feet. Which of the following statements are true?


 $\textbf{(A) }\text{Karl's garden is larger by 100 square feet.}$
$\textbf{(B) }\text{Karl's garden is larger by 25 square feet.}$
$\textbf{(C) }\text{The gardens are the same size.}$
$\textbf{(D) }\text{Makenna's garden is larger by 25 square feet.}$
$\textbf{(E) }\text{Makenna's garden is larger by 100 square feet.}$

 \textbf{Answer: }E



\textbf{Solution}

The area of a rectangle is given by the formula length times width. Karl's garden is $20 \times 45 = 900$ square feet and Makenna's garden is $25 \times 40 = 1000$ square feet. Since $1000 > 900,$ Makenna's garden is larger by $1000-900=100$ square feet. $\Rightarrow \boxed{ \textbf{(E)}\ \text{Makenna's garden is larger by 100 square feet.} }$



	\item (\textit{AMC 8 2011}) Extend the square pattern of 8 black and 17 white square tiles by attaching a border of black tiles around the square. What is the ratio of black tiles to white tiles in the extended pattern?


                    \begin{figure}[H]
                        \centering
                        \includegraphics[width=0.5\textwidth]{Problem/AMC_8_AMC_2011_3_Problem_Figure_1.png}
                    \end{figure}
                    
$\textbf{(A) }8:17 \qquad\textbf{(B) }25:49 \qquad\textbf{(C) }36:25 \qquad\textbf{(D) }32:17 \qquad\textbf{(E) }36:17$

 \textbf{Answer: }D



\textbf{Solution 1}

One way of approaching this is drawing the next circle of boxes around the current square.
                    \begin{figure}[H]
                        \centering
                        \includegraphics[width=0.5\textwidth]{Solution/AMC_8_AMC_2011_3_Solution_Figure_1.png}
                    \end{figure}
                    We can now count the number of black and white tiles; 32 black tiles and 17 white tiles. This means the answer is $\boxed{\textbf{(D) }32:17}$.

 



\textbf{Solution 2}

If we did not want to draw a diagram though that is probably simpler in this case, we can imagine the last border of black tiles. We see that each side length (if the side length of a square is $1$) increases by $2$ each time we go out one layer. Therefore, we know that the next border will have a side length of $7$. That border has $7^2-5^2=24$ black squares in it. The other black border has $3^2-1^2=8$ squares in it, so there are $32$ black squares in all. The other ones must all be white, so there are $49-32=17$ white squares. Thus, we get our answer of $\boxed{\textbf{(D) }32:17}$.

 



\textbf{Solution 3 (somewhat fake)}

Notice that in total, once we add the border there will be $7^2$ squares total. So far as there is only one option that sums to 49 and we know that there are an integer number of boxes, it has to be $\boxed{\textbf{(D) }32:17}$. ~math\_genius\_11

 


	\item (\textit{AMC 8 2011}) Here is a list of the numbers of fish that Tyler caught in nine outings last summer: 
 \[2,0,1,3,0,3,3,1,2.\] Which statement about the mean, median, and mode is true?


 $\textbf{(A) }\text{median} < \text{mean} < \text{mode} \qquad \textbf{(B) }\text{mean} < \text{mode} < \text{median} \\ \\ \textbf{(C) }\text{mean} < \text{median} < \text{mode} \qquad \textbf{(D) }\text{median} < \text{mode} < \text{mean} \\ \\ \textbf{(E) }\text{mode} < \text{median} < \text{mean}$

 \textbf{Answer: }C



\textbf{Solution}

First, put the numbers in increasing order.

 
 \[0,0,1,1,2,2,3,3,3\]

 The mean is $\frac{0+0+1+1+2+2+3+3+3}{9} = \frac{15}{9},$ the median is $2,$ and the mode is $3.$ Because, $\frac{15}{9} < 2 < 3,$ the answer is $\boxed{\textbf{(C)}\ \text{mean} < \text{median} < \text{mode}}$



	\item (\textit{AMC 8 2011}) What time was it $2011$ minutes after midnight on January 1, 2011?


 $\textbf{(A) }\text{January 1 at 9:31 PM}$
$\textbf{(B) }\text{January 1 at 11:51 PM}$
$\textbf{(C) }\text{January 2 at 3:11 AM}$
$\textbf{(D) }\text{January 2 at 9:31 AM}$
$\textbf{(E) }\text{January 2 at 6:01 PM}$

 \textbf{Answer: }D



\textbf{Solution 1}

There are $60$ minutes in an hour. $2011/60=33\text{r}31,$ or $33$ hours and $31$ minutes. There are $24$ hours in a day, so the time is $9$ hours and $31$ minutes after midnight on January 2, 2011. $\Rightarrow\boxed{\textbf{(D)}\text{ January 2 at 9:31AM}}$




\textbf{Solution 2}

Like the previous solution, $2011$ minutes is $33$ hours and $31$ minutes. That means it has to be January 1 at 9:31PM or January 2 at 9:31AM because they are the only ones that end with 31. But, the answer can’t be in January 1 because there are 24 hours in a day and 33 is greater than 24. So, the answer is $\Rightarrow\boxed{\textbf{(D)}\text{ January 2 at 9:31AM}}$.

 -RealCXY

 



\textbf{Solution 3 (Estimate)}

The answer choices are relatively far apart, so we can use estimation to solve this problem. $2100/60$ is around $35$ hours, or about $1.5$ days. The only answer choice around $1.5$ days from the midnight of January 1, 2011 is $\boxed{\textbf{(D)}\text{ January 2 at 9:31AM}}$.

 MegaBoy6679 :D 23:31, 26 December 2022 (EST)

 


	\item (\textit{AMC 8 2011}) In a town of 351 adults, every adult owns a car, motorcycle, or both. If 331 adults own cars and 45 adults own motorcycles, how many of the car owners do not own a motorcycle?


 $\textbf{(A) }20 \qquad\textbf{(B) }25 \qquad\textbf{(C) }45 \qquad\textbf{(D) }306 \qquad\textbf{(E) }351$

 \textbf{Answer: }D



\textbf{Solution 1}

By \href{https://artofproblemsolving.com/wiki/index.php?title=PIE}{PIE}, the number of adults who own both cars and motorcycles is $331+45-351=25.$ Out of the $331$ car owners, $25$ of them own motorcycles and $331-25=\boxed{\textbf{(D)}\ 306}$ of them don't.

 



\textbf{Solution 2}

There are $351$ total adults, and $45$ own a motorcycle. The number of adults that don't own a motorcycle is $351 - 45 = 306$. Since everyone owns a car or motorcycle and one who doesn't own a motorcycle owns a car, the answer is $\boxed{\textbf{(D)}\ 306}$.

 



\textbf{Solution 3 (answer choices)}

Clearly, we can eliminate answer choice $E$, for at least one adult owns motorcycles. It is also fairly obvious that the answer must be in the 300 range, giving us $\boxed{\textbf{(D)}\ 306}$



	\item (\textit{AMC 8 2011}) Each of the following four large congruent squares is subdivided into combinations of congruent triangles or rectangles and is partially bolded. What percent of the total area is partially bolded?


 
                    \begin{figure}[H]
                        \centering
                        \includegraphics[width=0.5\textwidth]{Problem/AMC_8_AMC_2011_7_Problem_Figure_1.png}
                    \end{figure}
                    
$\textbf{(A)}12\frac{1}{2}\qquad\textbf{(B)}20\qquad\textbf{(C)}25\qquad\textbf{(D)}33\frac{1}{3}\qquad\textbf{(E)}37\frac{1}{2}$

 \textbf{Answer: }C



\textbf{Solution}

Assume that the area of each square is $1$.  Then, the area of the bolded region in the top left square is $\dfrac{1}{4}$.  The area of the top right bolded region is $\dfrac{1}{8}$.  The area of the bottom left bolded region is $\dfrac{3}{8}$.  And the area of the bottom right bolded region is $\dfrac{1}{4}$.  Add the four fractions:  $\dfrac{1}{4} + \dfrac{1}{8} + \dfrac{3}{8} + \dfrac{1}{4} = 1$.  The four squares together have an area of $4$, so the percentage bolded is $\dfrac{1}{4} \cdot 100 = \boxed{\textbf{(C)}\ 25}$.

 


	\item (\textit{AMC 8 2011}) Bag A has three chips labeled 1, 3, and 5. Bag B has three chips labeled 2, 4, and 6. If one chip is drawn from each bag, how many different values are possible for the sum of the two numbers on the chips?


 $\textbf{(A) }4 \qquad\textbf{(B) }5 \qquad\textbf{(C) }6 \qquad\textbf{(D) }7 \qquad\textbf{(E) }9$

 \textbf{Answer: }B



\textbf{Solution 1}

By adding a number from Bag A and a number from Bag B together, the values we can get are $3, 5, 7, 5, 7, 9, 7, 9, 11.$ Therefore the number of different values is $\boxed{\textbf{(B)}\ 5}$.

 



\textbf{Solution 2}

The sum of an even number added to an odd number is always odd. The smallest possible sum is $3$, and the largest possible sum is $11$. The odd numbers in between can be achieved by replacing chips with ${\pm }2$ within the same bag. Therefore, we can conclude that there are $(11-3)/2(*)+1=\boxed{\textbf{(B)}\ 5}$ possible sums.

 
            \begin{itemize}
        \item Note that we are only accounting for odd numbers, thus the $/2$.
            \end{itemize}
            


	\item (\textit{AMC 8 2011}) Carmen takes a long bike ride on a hilly highway. The graph indicates the miles traveled during the time of her ride. What is Carmen's average speed for her entire ride in miles per hour?

                    \begin{figure}[H]
                        \centering
                        \includegraphics[width=0.5\textwidth]{Problem/AMC_8_AMC_2011_9_Problem_Figure_1.png}
                    \end{figure}
                    
$\textbf{(A) }2\qquad\textbf{(B) } 2.5\qquad\textbf{(C) } 4\qquad\textbf{(D) } 4.5\qquad\textbf{(E) } 5$

 \textbf{Answer: }E



\textbf{Solution}

We observe the graph and see that the shape of the graph does not matter. We only want the total time it took Carmen and the total distance she traveled. Based on the graph, Carmen traveled 35 miles for 7 hours. Therefore, her average speed is $\boxed{\textbf{(E)}\ 5}$



	\item (\textit{AMC 8 2011}) The taxi fare in Gotham City is $2.40 for the first $\frac12$ mile and additional mileage charged at the rate $0.20 for each additional 0.1 mile. You plan to give the driver a $2 tip. How many miles can you ride for $10?


 $\textbf{(A) } 3.0\qquad\textbf{(B) }3.25\qquad\textbf{(C) }3.3\qquad\textbf{(D) }3.5\qquad\textbf{(E) }3.75$

 \textbf{Answer: }C



\textbf{Solution 1 (Algebra)}

Let $x$ be the number of miles you ride. The number of miles you ride after the first half mile is $x-0.5.$ We can write this equation:

 


 
 \begin{align*} 10 &= 2.4 + 0.2 \times \frac{x-0.5}{0.1} + 2\\ 5.6 &= 2(x-0.5)\\ 2.8 &= x-0.5\\ x &= \boxed{\textbf{(C)}\ 3.3}\end{align*}





\textbf{Solution 2 (Table)}

\begin{array}{|c|c|c|}Miles & Money & Remark\\\hline0 & 8 & \text{2  for  tip} \\\hline0.5 & 5.6 & \text{-2.4  for  first  0.5  mi} \\\hline3.3 & 0 & \text{-5.6  for  2.8  mi}\end{array}

 Miles travelled = $\boxed{\textbf{(C)}\ 3.3}$

 Solution by Anshulb, formatted by Cheetahboy93

 


	\item (\textit{AMC 8 2011}) The graph shows the number of minutes studied by both Asha (black bar) and Sasha (grey bar) in one week. On the average, how many more minutes per day did Sasha study than Asha?

                    \begin{figure}[H]
                        \centering
                        \includegraphics[width=0.5\textwidth]{Problem/AMC_8_AMC_2011_11_Problem_Figure_1.png}
                    \end{figure}
                    
$\textbf{(A)}\ 6\qquad\textbf{(B)}\ 8\qquad\textbf{(C)}\ 9\qquad\textbf{(D)}\ 10\qquad\textbf{(E)}\ 12$

 \textbf{Answer: }A



\textbf{Solution}

Average the differences between each day. We get $10, -10,\text{ } 20,\text{ } 30,-20$. We find the average of this list to get $\boxed{\textbf{(A)}\ 6}$. ( In case you were wondering, the way to calculate the average is $\frac{(10+(-10)+20+30+(-20))}{5} = \frac{ 30}{5} = 6$. So the answer is indeed  $\boxed{\textbf{(A)}\ 6}$)

 



\textbf{Solution 2}

This solution may take longer to do than the first solution. In total, Asha studied for 400 minutes a week (80 minutes per day) and Sasha studied for 430 minutes a week (86 minutes per day).  86 - 80 = 6.  Therefore, the answer is $\boxed{\textbf{(A)}\ 6}$.

 -MiracleMaths

 


	\item (\textit{AMC 8 2011}) Angie, Bridget, Carlos, and Diego are seated at random around a square table, one person to a side. What is the probability that Angie and Carlos are seated opposite each other?


 $\textbf{(A) } \frac14 \qquad\textbf{(B) } \frac13 \qquad\textbf{(C) } \frac12 \qquad\textbf{(D) } \frac23 \qquad\textbf{(E) } \frac34$

 \textbf{Answer: }B



\textbf{Solution 1}

If we designate a person to be on a certain side, then all placements of the other people can be considered unique. WLOG, assign Angie to be on the side. There are then $3!=6$ total seating arrangements. If Carlos is across from Angie, there are only $2!=2$ ways to fill the remaining two seats. Then the probability Angie and Carlos are seated opposite each other is $\frac26=\boxed{\textbf{(B)}\ \frac13}$.

 



\textbf{Solution 2}

If we seat Angie first, there would be only one out of three ways Carlos can sit across from Angie. So the final answer is $\boxed{\textbf{(B) }\frac{1}{3}}$




\textbf{Solution 3 (cheese)}

Three people can be seated opposite of Angie (Bridget, Carlos, or Diego). Only one of these is Carlos. Thus, the answer is $\boxed{\textbf{(B) }\frac{1}{3}}$.- @scthecool

 


	\item (\textit{AMC 8 2011}) Two congruent squares, $ABCD$ and $PQRS$, have side length $15$. They overlap to form the $15$ by $25$ rectangle $AQRD$ shown. What percent of the area of rectangle $AQRD$ is shaded? 

                    \begin{figure}[H]
                        \centering
                        \includegraphics[width=0.5\textwidth]{Problem/AMC_8_AMC_2011_13_Problem_Figure_1.png}
                    \end{figure}
                    
$\textbf{(A)}\ 15\qquad\textbf{(B)}\ 18\qquad\textbf{(C)}\ 20\qquad\textbf{(D)}\ 24\qquad\textbf{(E)}\ 25$

 \textbf{Answer: }C



\textbf{Solution 2}

The length of BP is 5. the ratio of the areas is $\dfrac{5\cdot 15}{25\cdot 15}=\dfrac{5}{25}=20\%$-Megacleverstarfish15

 



\textbf{Solution 3(similar to Solution 2)}

To find the overlap length, we do the total length of the squares and subtract $25$(side length of figure). $(15 + 15) - 25 = 5$, so the overlap length is $5$. To find what percentage of $AQRD$ is shaded, we divide the shaded part by the area of the $AQRD$, so the percentage is $\dfrac{15 \cdot 5}{15 \cdot 25}$ = $\dfrac{5}{25}$ = $\dfrac{1}{5}$ = $\dfrac{20}{100}$ = $20$\%, so the answer is $\boxed{ \textbf{(C)}\ \text{20} }$.

 


	\item (\textit{AMC 8 2011}) There are $270$ students at Colfax Middle School, where the ratio of boys to girls is $5 : 4$. There are $180$ students at Winthrop Middle School, where the ratio of boys to girls is $4 : 5$. The two schools hold a dance and all students from both schools attend. What fraction of the students at the dance are girls?


 $\textbf{(A) } \dfrac7{18} \qquad\textbf{(B) } \dfrac7{15} \qquad\textbf{(C) } \dfrac{22}{45} \qquad\textbf{(D) } \dfrac12 \qquad\textbf{(E) } \dfrac{23}{45}$

 \textbf{Answer: }C



\textbf{Solution}

At Colfax Middle School, there are $\frac49 \times 270 = 120$ girls. At Winthrop Middle School, there are $\frac59 \times 180 = 100$ girls. The ratio of girls to the total number of students is $\frac{120+100}{270+180} = \frac{220}{450} = \boxed{\textbf{(C)}\ \frac{22}{45}}$




\textbf{Solution 2 (Guess and Check)}

Since we know there are $270$ kids at Colfax Middle School and we know the ratio of boys to girls is $5:4$, we can use guess and check to find the number of girls at the dance (and vice versa). $5:4=150:120$ and $4:5=80:100$. Now that the ratios sum to $270$ and $180$, we know that in Colfax there are $120$ girls and in Winthrop, there are $100$ girls. $120+100=220$ total girls and $270+180=450$ total people. $\frac{220}{450}=\boxed{\textbf{(C)} \frac{22}{45}}$




\textbf{Solution 3 (Educated Guess)}

We know that since $270$ is greater than $180$ and how the ratios are opposites, we know that there will be more boys than girls, so the fraction of students at the dance that are girls should be less than $\frac{1}{2}$. We also know that the amount of girls and the amount of boys are close, so the fraction of students at the dance that are girls should be really close to $\frac{1}{2}$. The only answer choice that/best satisfies these conditions is $\boxed{\textbf{(C)}\ \frac{22}{45}}$.

 


	\item (\textit{AMC 8 2011}) How many digits are in the product $4^5 \cdot 5^{10}$?


 $\textbf{(A) } 8 \qquad\textbf{(B) } 9 \qquad\textbf{(C) } 10 \qquad\textbf{(D) } 11 \qquad\textbf{(E) } 12$

 \textbf{Answer: }D



\textbf{Solution}


 \[4^5 \cdot 5^{10} = 2^{10} \cdot 5^{10} = 10^{10}.\]

 That is one $1$ followed by ten $0$'s, which is $\boxed{\textbf{(D)}\ 11}$ digits.

 



\textbf{Solution 2}

$4^5$ has $4$ digits while $5^{10}$ has $7$ digits. This means that $4^5 \cdot 5^{10}$ has a total of $7+4=\boxed{\textbf{(D)} 11}$ digits.

 


	\item (\textit{AMC 8 2011}) Let $A$ be the area of the triangle with sides of length $25, 25$, and $30$. Let $B$ be the area of the triangle with sides of length $25, 25,$ and $40$. What is the relationship between $A$ and $B$?


 $\textbf{(A) } A = \dfrac9{16}B \qquad\textbf{(B) } A = \dfrac34B \qquad\textbf{(C) } A=B \qquad\textbf{(D) } A = \dfrac43B \qquad \textbf{(E) }A = \dfrac{16}9B$

 \textbf{Answer: }C



\textbf{Solution 1 (Using the Pythagorean Theorem)}

\textbf{25-25-30}

 We can draw the altitude for the side with length 30. By HL Congruence, the two triangles formed are congruent. Thus the altitude splits the side with length 30 into two segments with length 15. By the \href{https://artofproblemsolving.com/wiki/index.php?title=Pythagorean\_Theorem}{Pythagorean Theorem}, we have 
 \[15^2 + x^2 =25^2\]
 \[x^2 = 25^2 - 15^2\]
 \[x^2 = (25 + 15)(25-15)\]
 \[x^2= 40\cdot 10\]
 \[x^2= 400\]
 \[x = \sqrt{400}\]
 \[x= 20\]

 
Thus we have two 15-20-25 right triangles.

 \textbf{25-25-40}

 We can draw the altitude for the side with length 40. By HL Congruence, the two triangles formed are congruent. Thus the altitude splits the side with length 40 into two segments with length 20. From the 25-25-30 case, we know that the other side length is 15, so we have two 15-20-25 right triangles.Let the area of a 15-20-25 right triangle be $x$.
 \[a = 2x\]
 \[b = 2x\]
 \[\boxed{\textbf{(C) } A = B}\]




\textbf{Solution 2 (Using Heron's Formula)}

Using Heron's formula, we can calculate the area of the two triangles. The formula states that
 \[A = \sqrt{s(s - a)(s - b)(s - c)}\]where $s$ is the semiperimeter of a triangle with side lengths $a$, $b$, and $c$.

 For the 25-25-30 triangle, 
 \[s = \frac{25 + 25 + 30}{2} = 40\]Therefore,
 \[A = \sqrt{40 \cdot 15 \cdot 15 \cdot 10} = 300\]

 For the 25-25-40 triangle, 
 \[s = \frac{25 + 25 + 40}{2} = 45\]Therefore,
 \[B = \sqrt{45 \cdot 20 \cdot 20 \cdot 5} = 300\]

 Hence, 
 \[A = B \hspace{0.15in} \Longrightarrow \hspace{0.15in} \boxed{\textbf{(C)}}\]




\textbf{Solution 3 (Simple)}

\textbf{ Assume } the triangles are right. Then, the 25-25-30 has hypotenuse 30, whereas 25-25-40 has hypotenuse 40. Because the formula for a triangles' area is $(bh)/2$, the hypotenuse does not matter. So, since the base and height are the same, the area will also be the same meaning 
 \[A = B \hspace{0.15in} \Longrightarrow \hspace{0.15in} \boxed{\textbf{(C)}}\]



	\item (\textit{AMC 8 2011}) Let $w$, $x$, $y$, and $z$ be whole numbers. If $2^w \cdot 3^x \cdot 5^y \cdot 7^z = 588$, then what does $2w + 3x + 5y + 7z$ equal?


 $\textbf{(A) } 21\qquad\textbf{(B) }25\qquad\textbf{(C) }27\qquad\textbf{(D) }35\qquad\textbf{(E) }56$

 \textbf{Answer: }A



\textbf{Solution}

The \href{https://artofproblemsolving.com/wiki/index.php?title=Prime\_factorization}{prime factorization} of $588$ is $2^2\cdot3\cdot7^2.$ We can see $w=2, x=1,$ and $z=2.$ Because $5^0=1, y=0.$

 \[2w+3x+5y+7z=4+3+0+14=\boxed{\textbf{(A)}\ 21}\]



	\item (\textit{AMC 8 2011}) A fair 6-sided die is rolled twice. What is the probability that the first number that comes up is greater than or equal to the second number?


 $\textbf{(A) }\dfrac16\qquad\textbf{(B) }\dfrac5{12}\qquad\textbf{(C) }\dfrac12\qquad\textbf{(D) }\dfrac7{12}\qquad\textbf{(E) }\dfrac56$

 \textbf{Answer: }D



\textbf{Solution}

There are $6\cdot6=36$ ways to roll the two dice, and 6 of them result in two of the same number. Out of the remaining $36-6=30$ ways, the number of rolls where the first dice is greater than the second should be the same as the number of rolls where the second dice is greater than the first. In other words, there are $30/2=15$ ways the first roll can be greater than the second. The probability the first number is greater than or equal to the second number is

 
 \[\frac{15+6}{36}=\frac{21}{36}=\boxed{\textbf{(D)}\ \frac{7}{12}}\]




\textbf{Solution 2}

If we list out some ways the first dice is greater than the second dice, we might see a pattern.

 If dice $2$ is $1$, then dice $1$ can be $1, 2, 3, 4, 5,$ or $6$.

 If dice $2$ is $2$, then dice $1$ can be $2, 3, 4, 5,$ or $6$.

 If dice $2$ is $3$, then dice $1$ can be $3, 4, 5,$ or $6$.

 In the first case, dice $1$ had six values it could land on. In the second case, dice $1$ had five values it could land on. In the third case, dice $1$ had four values it could land on. If we keep going, we would eventually get the pattern $6+5+4+3+2+1=21$ as the number of total values dice $1$ could land on if it's greater or equal to dice $2$. Since there are $36$ ways two dice could land, our answer is $\frac{21}{36}$, or $\boxed{\textbf{(D)} \frac{7}{12}}$.

 


	\item (\textit{AMC 8 2011}) How many rectangles are in this figure?


 
                    \begin{figure}[H]
                        \centering
                        \includegraphics[width=0.5\textwidth]{Problem/AMC_8_AMC_2011_19_Problem_Figure_1.png}
                    \end{figure}
                    
$\textbf{(A)}\ 8\qquad\textbf{(B)}\ 9\qquad\textbf{(C)}\ 10\qquad\textbf{(D)}\ 11\qquad\textbf{(E)}\ 12$

 \textbf{Answer: }D



\textbf{Solution 1 (Splitting It Up)}

The figure can be divided into $7$ sections. The number of rectangles with just one section is $3.$ The number of rectangles with two sections is $5.$ There are none with only three sections. The number of rectangles with four sections is $3.$
$3+5+3=\boxed{\textbf{(D)}\ 11}$




\textbf{Solution 2 (Count by Reduction)}

We can remove the 3 big blocks of rectangles one by one. 7 (left) + 3 (bottom) + 1 = 11 are removed in total. 

 


	\item (\textit{AMC 8 2011}) Quadrilateral $ABCD$ is a trapezoid, $AD = 15$, $AB = 50$, $BC = 20$, and the altitude is $12$. What is the area of the trapezoid?


 
                    \begin{figure}[H]
                        \centering
                        \includegraphics[width=0.5\textwidth]{Problem/AMC_8_AMC_2011_20_Problem_Figure_1.png}
                    \end{figure}
                    
$\textbf{(A) }600\qquad\textbf{(B) }650\qquad\textbf{(C) }700\qquad\textbf{(D) }750\qquad\textbf{(E) }800$

 \textbf{Answer: }D



\textbf{Solution}


                    \begin{figure}[H]
                        \centering
                        \includegraphics[width=0.5\textwidth]{Solution/AMC_8_AMC_2011_20_Solution_Figure_1.png}
                    \end{figure}
                    

 If you draw altitudes from $A$ and $B$ to $CD,$ the trapezoid will be divided into two right triangles and a rectangle. You can find the values of $a$ and $b$ with the \href{https://artofproblemsolving.com/wiki/index.php?title=Pythagorean\_theorem}{Pythagorean theorem}.

 
 \[a=\sqrt{15^2-12^2}=\sqrt{81}=9\]

 \[b=\sqrt{20^2-12^2}=\sqrt{256}=16\]
$ABYX$ is a rectangle so $XY=AB=50.$

 \[CD=a+XY+b=9+50+16=75\]

 The area of the trapezoid is

 
 \[12\cdot \frac{(50+75)}{2} = 6(125) = \boxed{\textbf{(D)}\ 750}\]



	\item (\textit{AMC 8 2011}) Students guess that Norb's age is $24, 28, 30, 32, 36, 38, 41, 44, 47$, and $49$. Norb says, "At least half of you guessed too low, two of you are off by one, and my age is a prime number." How old is Norb?


 $\textbf{(A) }29\qquad\textbf{(B) }31\qquad\textbf{(C) }37\qquad\textbf{(D) }43\qquad\textbf{(E) }48$

 \textbf{Answer: }C



\textbf{Solution 1}

At least half the guesses are too low, so Norb's age must be greater than $36$.

 If two of the guesses are off by one, then his age is in between two guesses whose difference is $2$. It could be $31,37,$ or $48,$ but because his age is greater than $36$ it can only be $37$ or $48$.

 Lastly, Norb's age is a prime number so the answer must be $\boxed{\textbf{(C)}\ 37}$.

 



\textbf{Solution 2 (Alternative approach)}

Since two guesses are off by one, we know that both $x+1$ and $x-1$ are in the list where $x$ is the age of Norb. Now, we know that $x+1$ and $x-1$ are $28$ and $30$, $30$ and $32$, $36$ and $38$ and $47$ and $49$. From these values, we know that $x$ must be $29$, $31$, and $37$. Since half of the guesses are too low, $24, 28, 30, 32,$ and $36$ are all too low so we can eliminate all numbers in our list lesser than or equal to $36$. Therefore, our list has only $37$ left so the answer is $\boxed{\textbf{(C)}\ 37}$.

 



\textbf{Solution 3 (Using the Answer Choices)}

If his age is $29$, then there won't be half of them who guessed less than his age. If his age is $31$, then there still won't be half of them who guessed less than his age. If his age is $37$, then there are half of them who guessed too low, half of them who guessed too high, and there are two people who guessed off by one. So, our answer is probably $\textbf{(C)}\ 37$, but let's still check the other two choices. If his age is $43$, then there won't be two people who were off by one. Finally, if his age is $48$, then his age won't be a prime number. Therefore, our answer is $\boxed{\textbf{(C)}\ 37}$.

 


	\item (\textit{AMC 8 2011}) What is the \textbf{tens} digit of $7^{2011}$?


 $\textbf{(A) }0\qquad\textbf{(B) }1\qquad\textbf{(C) }3\qquad\textbf{(D) }4\qquad\textbf{(E) }7$

 \textbf{Answer: }D



\textbf{Solution 1}

Since we want the tens digit, we can find the last two digits of $7^{2011}$.  We can do this by using modular arithmetic.
 \[7^1\equiv 07 \pmod{100}.\]
 \[7^2\equiv 49 \pmod{100}.\]
 \[7^3\equiv 43 \pmod{100}.\]
 \[7^4\equiv 01 \pmod{100}.\]We can write $7^{2011}$ as $(7^4)^{502}\times 7^3$.  Using this, we can say:
 \[7^{2011}\equiv (7^4)^{502}\times 7^3\equiv 7^3\equiv 343\equiv 43\pmod{100}.\]From the above, we can conclude that the last two digits of $7^{2011}$ are 43.  Since they have asked us to find the tens digit, our answer is $\boxed{\textbf{(D)}\ 4}$.

 -Ilovefruits

 



\textbf{Solution 2}

The last two digits of $7^{2011}$ can be found by taking $\pmod{100}$ of $7^{2011}$. By \href{https://artofproblemsolving.com/wiki/index.php?title=Euler\%27s\_Totient\_Theorem}{Euler's Totient Theorem}, we have that $7^{\phi(100)}\equiv 7^{40}\equiv 1 \pmod{100}$, so we can conclude that $7^{2011} \equiv 7^{40^{50}} \cdot 7^{11}\equiv 7^{11} \pmod{100}$. We can find that the modular inverse of $7 \pmod{100}$ is $43$, and since $7^{12}\equiv 01 \pmod{100}$, we can conclude that $7^{11}\equiv 7^{12} \cdot 7^{-1}\equiv 1 \cdot 43\equiv 43 \pmod{100}$. The problem asked for the tens digit, so our answer is $\boxed{\textbf{(D)}\ 4}$.~Baihly2024

 


	\item (\textit{AMC 8 2011}) How many 4-digit positive integers have four different digits, where the leading digit is not zero, the integer is a multiple of 5, and 5 is the largest digit?


 $\textbf{(A) }24\qquad\textbf{(B) }48\qquad\textbf{(C) }60\qquad\textbf{(D) }84\qquad\textbf{(E) }108$

 \textbf{Answer: }D



\textbf{Solution 1}

We can separate this into two cases. If an integer is a multiple of $5,$ the last digit must be either $0$ or $5.$

 Case 1: The last digit is $5.$ The leading digit can be $1,2,3,$ or $4.$ Because the second digit can be $0$ but not the leading digit, there are also $4$ choices. The third digit cannot be the leading digit or the second digit, so there are $3$ choices. The number of integers is this case is $4\cdot4\cdot3\cdot1=48.$

 Case 2: The last digit is $0.$ Because $5$ is the largest digit, one of the remaining three digits must be $5.$ There are $3$ ways to choose which digit should be $5.$ The remaining digits can be $1,2,3,$ or $4,$ but since they have to be different there are $4\cdot3$ ways to choose. The number of integers in this case is $1\cdot3\cdot4\cdot3=36.$

 Therefore, the answer is $48+36=\boxed{\textbf{(D)}\ 84}$.

 


	\item (\textit{AMC 8 2011}) In how many ways can 10001 be written as the sum of two primes?


 $\textbf{(A) }0\qquad\textbf{(B) }1\qquad\textbf{(C) }2\qquad\textbf{(D) }3\qquad\textbf{(E) }4$

 \textbf{Answer: }A



\textbf{Solution}

For the sum of two numbers to be odd, one must be odd and the other must be even, because all odd numbers are of the form $2n+1$ where n is an integer, and all even numbers are of the form $2m$ where m is an integer.
 \[2n + 1 + 2m = 2m + 2n + 1 = 2(m+n) + 1\] and $m+n$ is an integer because $m$ and $n$ are both integers.The only even prime number is $2,$ so our only combination could be $2$ and $9999.$ But, $9999$ is clearly divisible by $3$,  so the number of ways $10001$ can be written as the sum of two primes is $\boxed{\textbf{(A)}\ 0}$.

 



\textbf{Solution 2(Simple)}

First, we noticed that 10001 is equal to 5000+5001, if you subtract n to 5000 and add n to 5001, you always get an even number, even number is never a prime number except 2. We also see that whenever an addend is an odd number, the other addend will be even, so having an odd number as an addend is not possible, other than 9999 and 2, because 2 is a prime. We try 2 and 9999 but we can see 9999 is divisible by 3 and 9 clearly. So the answer is $\boxed{\textbf{(A)} 0}$




\textbf{Solution 3 (assumed previous knowledge)}

It is helpful to know and understand the Goldbach Conjecture - that every even number can be written as the sum of $2$ primes - and also, that the ${\textbf{odd numbers}\ }$ that are the sum of two primes are exactly two more than a prime. This is because to make the sum of two numbers odd, you must have one even and one odd. There is only one even prime, which is two, so the sum will be of the form $2+p$. Hence, the odd numbers that are the sum of two primes are exactly $2$ more than a prime. Relating to the problem, $10001$ is not $2$ more than a prime, because $10001-2=9999$ and we can easily see that $9999$ is divisible by $3$. Therefore, $10001$ cannot be written as the sum of two primes, and so therefore the answer is $\boxed{\textbf{(A)}\ 0}$~mk

 Note: the Goldbach Conjecture isn't solved, however you can assume it to be true for basically all numbers you see in competetions because many supercomputers tried to find a counterexample (went really high) but they didn't.

 For proof based competitions where the graders actually like read your solutions don't use it they won't accept the conjecture as a solution

 



\textbf{Solution 4 (similar to solution 1)}

Our equation can be rewritten to $Odd = Even + Odd$. Since both addends have to be prime, one has to be even and one prime, which means it's $2$. Plugging this in, the other number has to be $9999$ which is not prime (divisible by $9$) meaning there are $\boxed{\textbf{(A)}\ 0}$ answers. 

 


	\item (\textit{AMC 8 2011}) A circle with radius $1$ is inscribed in a square and circumscribed about another square as shown. Which fraction is closest to the ratio of the circle's shaded area to the area between the two squares?


 
                    \begin{figure}[H]
                        \centering
                        \includegraphics[width=0.5\textwidth]{Problem/AMC_8_AMC_2011_25_Problem_Figure_1.png}
                    \end{figure}
                    
$\textbf{(A)}\ \frac{1}2\qquad\textbf{(B)}\ 1\qquad\textbf{(C)}\ \frac{3}2\qquad\textbf{(D)}\ 2\qquad\textbf{(E)}\ \frac{5}2$

 \textbf{Answer: }A



\textbf{Solution 1}

The area of the smaller square is one half of the product of its diagonals. Note that the distance from a corner of the smaller square to the center is equivalent to the circle's radius so the diagonal is equal to the diameter: $2 \cdot  2 \cdot \frac{1}{2}=2.$

 The circle's shaded area is the area of the smaller square subtracted from the area of the circle: $\pi - 2.$

 If you draw the diagonals of the smaller square, you will see that the larger square is split $4$ congruent half-shaded squares. The area between the squares is equal to the area of the smaller square: $2.$

 Approximating $\pi$ to $3.14,$ the ratio of the circle's shaded area to the area between the two squares is about

 
 \[\frac{\pi-2}{2} \approx \frac{3.14-2}{2} = \frac{1.14}{2} \approx \boxed{\textbf{(A)}\ \frac12}\].

 



\textbf{Solution 2}

For the ratio of the circle's shaded area to the area between the squares to be $1,$ they would have to be approximately the same size. For any ratio larger than that, the circle's shaded area must be greater. However, we can clearly see that the circle's shaded area is part of the area between the squares, and is approximately $\boxed{\textbf{(A)}\ \frac12}$. The ratio must be less than 1, because the area inside the circle is less than the total shaded area, and the only answer choice that fits that is ${\frac12}$.Note that this solution is not rigorous, because we still should show that the ratio is less than $\frac{3}{4}$.

 



\textbf{Solution 3}

Set the side length of the bigger square to be $8$. Then the area of the big square is $8^2 =64$ and the area of the small square $(4\sqrt{2})^2  = 32$, making the difference $32$. The area of the circle is $4^2$ times $\pi$ which is $16 \pi$ or about $48$. Knowing the area of the small square is $32$. $48-32$ is $16$. The area of the big square is $64$. So $32/64$ is $1/2$, or $\boxed{\textbf{(A)}\ \frac12}$.

 



\textbf{Solution 4 (Non-rigorous speed solution)}

Note that that the shaded area of the circle is existent only between the gap between the two squares. This implies that the shaded area is less than the gap, and the only answer that shows that is $\boxed{\textbf{(A)}\ \frac{1}{2}}$.

 
- by nickelslordm

 


\end{enumerate}